% !PRINTED 25
\weberpage{25}

%<<CONT
% the German page break falls inside "das auf $Q$ | folgende Glied"; the English absorbs
% "following" onto the previous page, so this page opens directly with $Q_1$
$Q_1$, is not a normal divisor of $P$. Then among the groups $a^{-1} Q_1 a$ conjugate with respect to $P$ there are several distinct ones, which we shall denote by

\weq{1}{Q_1,\ Q_1',\ Q_1'',\ Q_1''' \cdots}

\gesperrt{The greatest common divisor $R$ of all these groups is again a normal divisor of $P$.}

Now all the groups (1) are greatest normal divisors of $Q$, as follows by the concluding theorem of \S.\,3 from the fact that $Q, Q_1$ are isomorphic with $a^{-1} Q a, a^{-1} Q_1 a$, while $Q$, as a normal divisor of $P$, is identical with $a^{-1} Q a$. We can therefore let any one of the groups (1), which all have the same index $\nu$, follow $Q$ in the composition series.

If now $Q_2$ is the greatest common divisor of $Q_1, Q_1'$ (and hence distinct from $Q_1$), then by the theorem \S.\,6, 4. $Q_2$ is a greatest normal divisor of $Q_1$ and $Q_1'$, whose index is likewise $\nu$. We can therefore continue the composition series from $Q$ on in the following two ways:

\weq{2}{
\begin{array}{ccc@{\qquad}ccc}
Q, & Q_1, & Q_2, & Q & Q_1' & Q_2 \\
& \nu & \nu & & \nu & \nu
\end{array}
}

We now wish to generalise the law expressed herein and to prove it by complete induction.

If $Q_2$ is not yet equal to $R$, we adjoin to $Q_1, Q_1'$ a third group $Q_1''$ from the sequence (1), so that the intersection $Q_3$ of $Q_1, Q_1', Q_1''$ becomes a proper divisor of $Q_2$, and we continue in this way until we have arrived at the greatest common divisor $R$ of all the groups (1).

Accordingly, let:

\weq{3}{
\begin{array}{l}
Q_2 \text{ the intersection of } Q_1,\ Q_1' \\
Q_3 \text{ \ ,,\qquad ,,\qquad ,,\ } Q_1,\ Q_1',\ Q_1'' \\
\cdot\ \cdot\ \cdot\ \cdot\ \cdot\ \cdot\ \cdot\ \cdot\ \cdot\ \cdot\ \cdot\ \cdot\ \cdot \\
Q_r \text{ \ ,,\qquad ,,\qquad ,,\ } Q_1,\ Q_1',\ Q_1'' \cdots Q_1^{(r-1)}
\end{array}
}

and let the arrangement of the groups $Q_1, Q_1' \cdots Q_1^{(r-1)}$ be so chosen that, of the groups $Q_1, Q_2 \cdots Q_r$, each is a proper divisor of the one preceding it. This arrangement can be continued so long as $Q_r$ is not equal to $R$.

We wish to prove that we can choose the composition series of $P$ in such a way that in it

\weq{4}{
\begin{array}{cccc}
Q, & Q_1, & Q_2 & \cdots\ Q_r\\
& \nu' & \nu & \nu % [sic: the source prints \nu' although the text says all indices are \nu.]
\end{array}
}
%CONT>>
